\subsection{Next steps}\label{sec:next}

The prototype turns a translated specification into review questions with solver evidence. The order example gives the author R2 and R3, a feasible no-stock submission, and incompatible obligations. For deletion, the reviewer also needs the assumed relation between removal and retention. Assumptions and alternative readings remain part of the review.

The experiments support continuing this approach with qualifications. Outcome-pair questions and a second reader improve detection while adding review work. Change checking now attributes tested removal gaps and offers saved-base reuse. E20 shows fewer requests alongside lower conflict detection, so full formalization remains the default. Input, runtime and representation limits remain; production accuracy and human review savings have not been established.

\paragraph{Independent evaluation.} Prepare a fresh set before another design change; the current sets have informed development. Add independently authored requirements and a real OpenSpec project's history. Report distinct defects as well as repeated runs, retain raw and doubtful-filtered findings, and measure review time. Record the requirements, reviewer decisions and fixed configuration so others can assess the results.

\paragraph{Change analysis on larger projects.} Test sentence-only edits, requirements without scenarios, state-only removals, and changes to old vocabulary meanings. Fix or reject unsupported edits explicitly. Compare repeated alternatives against one base with rebuilding after each archive; test cache validity across sampling, prompt and model changes. Investigate the E20 misses on fresh cases and measure context limits on a real project. The order removal remains a useful regression example: submission stays in scope while its no-stock branch is removed.

\paragraph{Vocabulary and representation.} Use the cloakroom failure and payroll miss as design cases: express retention as an obligation, test the outcome-pair cap, and investigate a doubt marker for collapsed lifecycles. Preserve contradictions such as order backordering while testing changes on fresh cases. The model remains one snapshot; a temporal extension has not been agreed. Reviewers still need to see where time and separate objects have been collapsed.

The order of work is evaluation, change analysis, then vocabulary selection. These are research priorities, without a staffing estimate or delivery schedule.
\source{plan.md, sections 2 and 5; research.md, sections 4 and 5}

